% ======================================================================
\chapter[Слабое взаимодействие. \texorpdfstring{$W^\pm$, $Z^0$}{W±, Z0} бозоны. Слабые распады]%
        {Слабое взаимодействие. \texorpdfstring{$W^\pm$, $Z^0$}{W±, Z0} бозоны. Слабые распады}
\label{t:16}
% ======================================================================

\begin{qbox}
\textbf{Билет 16.} Слабое взаимодействие. $W^\pm$, $Z^0$ бозоны, слабые распады лептонов и адронов.
\end{qbox}

\section{Содержание билета}

\textbf{Подтемы:}
\begin{enumerate}
 \item общие свойства слабого взаимодействия: универсальность, малый радиус, изменение аромата, нарушение чётности;
 \item переносчики $W^\pm$ и $Z^0$: заряженные и нейтральные токи, массы, открытие; теория Ферми и её связь с массой $W$;
 \item электрослабое объединение (качественно);
 \item слабые распады лептонов: распад мюона и $\tau$-лептона, закон Сарджента;
 \item слабые распады адронов: $\beta$-распад нейтрона, распады $\pi$, $K$, $\Lambda$; угол Кабиббо;
 \item законы сохранения и их нарушения в слабых процессах.
\end{enumerate}

\textbf{Необходимые понятия:} Стандартная модель и обменный механизм (\tref{t:14}); кварковый состав
адронов (\tref{t:15}); спиральность; время жизни и ширина (\tref{t:02}).

\section{Определения и физические принципы}

\begin{definition}[слабое взаимодействие]
Фундаментальное взаимодействие, в котором участвуют все лептоны и кварки; единственное, способное
изменять аромат фермионов ($d\to u$, $\mu\to\nu_\mu$, $s\to u$). Переносчики — массивные векторные бозоны $W^+$, $W^-$
($M_W=80{,}4$~ГэВ) и $Z^0$ ($M_Z=91{,}2$~ГэВ).
\end{definition}

\begin{definition}[заряженные и нейтральные токи]
\emph{Заряженный ток} — процесс с обменом $W^\pm$: заряд лептона или кварка меняется на единицу ($e^-\leftrightarrow\nu_e$,
$d\leftrightarrow u$). \emph{Нейтральный ток} — обмен $Z^0$: фермион остаётся тем же ($\nu_\mu e^-\to\nu_\mu e^-$).
Нейтральные токи с изменением аромата (например, $s\to d$ через $Z^0$) в низшем порядке отсутствуют.
\end{definition}

\begin{definition}[спиральность]
Проекция спина на направление импульса, $h=\vect s\cdot\vect p/|\vect p|$ (в единицах $\hbar$). Левая спиральность — $h=-\frac12$.
В заряженных токах участвуют только левые (по киральности) фермионы и правые антифермионы; для безмассовых
частиц это совпадает со спиральностью: нейтрино левые, антинейтрино правые.
\end{definition}

\emph{Физический смысл.} Слабое взаимодействие отвечает за процессы, недоступные сильному и
электромагнитному: бета-распад, распады мюона, пионов, странных частиц, термоядерный $pp$-цикл на Солнце
($p+p\to d+e^++\nu_e$), рассеяние нейтрино.

\emph{Область применимости описания.} При энергиях $E\ll M_Wc^2$ взаимодействие эффективно точечное (теория
Ферми с константой $G_F$). При $E\sim M_Wc^2$ необходимо явное описание обмена $W$, $Z$ в рамках электрослабой теории.

\section{Основная теория}

\subsection{Проблема и история}

Бета-распад с непрерывным спектром электронов привёл к гипотезе нейтрино (Паули, 1930) и к первой теории
слабого взаимодействия (Ферми, 1934): четыре фермиона взаимодействуют в одной точке с константой
$G_F$. Теория описывала распады, но предсказывала неограниченный рост сечений с энергией ($\sigma\propto G_F^2E^2$),
что недопустимо физически (нарушение унитарности при $E\sim300$~ГэВ). Решение — переносчик конечной массы:
при $E\ll Mc^2$ взаимодействие выглядит точечным, а при больших энергиях рост сечений прекращается.

\subsection{Свойства слабого взаимодействия}

\begin{enumerate}
 \item \textbf{Малый радиус:} $r_W=\hbar/M_Wc\approx2{,}5\cdot10^{-3}$~фм (\tref{t:14}).
 \item \textbf{«Слабость» при низких энергиях:} амплитуда обмена $W$ пропорциональна $g^2/(q^2-M_W^2c^2)$; при
       $q^2\ll M_W^2c^2$ это константа $\propto g^2/M_W^2$, т.\,е. точечное взаимодействие с
       \begin{equation}\label{eq:16:GF}
         \frac{G_F}{\sqrt2}=\frac{g^2}{8M_W^2}\quad(\hbar=c=1),\qquad \frac{G_F}{(\hbar c)^3}=1{,}166\cdot10^{-5}\ \text{ГэВ}^{-2}.
       \end{equation}
       Сама константа $g$ сравнима с электрическим зарядом ($e=g\sin\theta_W$, $\sin^2\theta_W\approx0{,}23$).
 \item \textbf{Универсальность:} лептоны и кварки (с учётом смешивания) взаимодействуют с $W$ с одной константой.
 \item \textbf{Изменение аромата} — только через $W^\pm$; поэтому лёгчайшие частицы каждого аромата
       ($\mu$, $\pi^\pm$, $K$, $\Lambda$, нейтрон) распадаются только слабо и живут долго.
 \item \textbf{Нарушение чётности ($P$) и зарядового сопряжения ($C$)}, слабое нарушение $CP$.
\end{enumerate}

\paragraph{Нарушение чётности.} Т.~Ли и Ч.~Янг (1956) указали, что сохранение чётности в слабых процессах не
проверено. В опыте Ц.~Ву (1957) ядра $^{60}$Co поляризовались магнитным полем при $T\approx0{,}01$~К; электроны
бета-распада вылетали преимущественно \emph{против} спина ядра. При зеркальном отражении импульс (полярный вектор) меняет
знак, а спин (аксиальный вектор) — нет, поэтому зеркальный процесс имел бы преимущественный вылет \emph{по} спину.
Асимметрия $\avg{\vect J\cdot\vect p_e}\neq0$ доказывает, что слабое взаимодействие различает правое и левое. Опыт Гольдхабера
(1958) показал, что нейтрино имеет левую спиральность.

\subsection{Переносчики \texorpdfstring{$W^\pm$ и $Z^0$}{W± и Z0}. Электрослабое объединение}

В электрослабой теории (Глэшоу, Вайнберг, Салам; Нобелевская премия 1979) электромагнитное и слабое
взаимодействия — проявления единого взаимодействия с группой $SU(2)_L\times U(1)_Y$. Четыре калибровочных бозона
после нарушения симметрии механизмом Хиггса (\tref{t:14}) становятся $W^\pm$, $Z^0$ (массивные) и фотоном (безмассовым).
Связь масс и констант:
\begin{equation}\label{eq:16:ew}
  e=g\sin\theta_W,\qquad M_Z=\frac{M_W}{\cos\theta_W},\qquad
  M_W^2=\frac{\sqrt2\,g^2}{8G_F}=\frac{\pi\alpha}{\sqrt2G_F\sin^2\theta_W}.
\end{equation}
Подстановка $\alpha$, $G_F$ и $\sin^2\theta_W\approx0{,}23$ (измеренного по нейтральным токам) даёт $M_W\approx78$~ГэВ (без учёта радиационных
поправок) — предсказание, подтверждённое открытием $W$ и $Z$ в ЦЕРН (1983) на $p\bar p$-коллайдере: $W$ — по
распадам $W\to e\nu$, $Z$ — по парам $e^+e^-$, $\mu^+\mu^-$. Нейтральные токи впервые наблюдались в 1973~г. (камера Gargamelle,
рассеяние $\bar\nu_\mu e^-\to\bar\nu_\mu e^-$ и $\nu N\to\nu X$).

Свойства: $W$ распадается на $\ell\nu$ (по $\approx11\,\%$ на каждый лептон) и $q\bar q'$; полная ширина $\Gamma_W\approx2{,}1$~ГэВ,
время жизни $\sim3\cdot10^{-25}$~с. Ширина $Z^0$ ($\Gamma_Z=2{,}50$~ГэВ) включает распады в $\nu\bar\nu$; её измерение на LEP показало,
что существует ровно три типа лёгких нейтрино.

\subsection{Слабые распады лептонов}

\paragraph{Распад мюона:}
\begin{equation}\label{eq:16:mu}
  \mu^-\to e^-+\bar\nu_e+\nu_\mu,\qquad \tau_\mu=2{,}197\ \text{мкс}.
\end{equation}
Сохраняются заряд, $L_\mu$ ($1=1$) и $L_e$ ($0=1-1$). Процесс — чисто лептонный, без сильных эффектов, поэтому вычисляется
точно (без вывода):
\begin{equation}\label{eq:16:Gmu}
  \Gamma_\mu=\frac{\hbar}{\tau_\mu}=\frac{G_F^2(m_\mu c^2)^5}{192\pi^3(\hbar c)^6}.
\end{equation}
Именно из $\tau_\mu$ определяется $G_F$. Зависимость $\Gamma\propto G_F^2m^5$ («закон Сарджента») получается из соображений
размерности (вывод~\ref{der:16:sargent}).

\paragraph{Распады $\tau$-лептона.} $\tau^-$ тяжелее и может распадаться на $e\bar\nu_e\nu_\tau$, $\mu\bar\nu_\mu\nu_\tau$ и адроны+$\nu_\tau$.
Проверка универсальности (вывод~\ref{der:16:tau}): по $\tau_\mu$, отношению масс и доле распада $\tau\to e\bar\nu\nu$ ($\approx17{,}8\,\%$)
предсказывается $\tau_\tau\approx2{,}9\cdot10^{-13}$~с — в согласии с опытом.

\subsection{Слабые распады адронов}

На кварковом уровне все заряженно-токовые распады адронов — превращения $d\to u+e^-+\bar\nu_e$, $s\to u+\dots$ через
виртуальный $W^-$ (рис.~\ref{fig:16:decays}).

\begin{enumerate}
 \item \textbf{Бета-распад нейтрона:} $n\to p+e^-+\bar\nu_e$ ($udd\to uud$, т.\,е. $d\to ue^-\bar\nu_e$). $Q=0{,}782$~МэВ, $\tau_n\approx880$~с.
       Огромное время жизни при той же константе — следствие малого энерговыделения ($\Gamma\propto Q^5$).
 \item \textbf{Распад пиона:} $\pi^+\to\mu^++\nu_\mu$ ($u\bar d\to W^+\to\mu^+\nu_\mu$), $\tau=2{,}6\cdot10^{-8}$~с. Канал $\pi^+\to e^+\nu_e$
       энергетически выгоднее, но подавлен в $\approx8\cdot10^3$ раз: пион имеет спин 0, поэтому заряженный лептон и нейтрино
       вылетают с противоположными импульсами и одинаковыми спиральностями; нейтрино левое, значит $\ell^+$ тоже
       должен быть левым, что для антилептона возможно лишь с амплитудой $\propto m_\ell/E$. Подавление $\propto m_e^2$ —
       прямое проявление $V\!-\!A$ структуры (только левые токи).
 \item \textbf{Распад $\pi^0\to2\gamma$} ($\tau=8\cdot10^{-17}$~с) — электромагнитный, для сравнения.
 \item \textbf{Распады странных частиц:} $K^+\to\mu^+\nu_\mu$ ($\tau=1{,}24\cdot10^{-8}$~с), $\Lambda\to p+\pi^-$ ($\tau=2{,}6\cdot10^{-10}$~с) —
       $s\to u$, странность меняется на $\Delta S=1$. Сильное взаимодействие сохраняет странность, поэтому странные частицы
       рождаются парами быстро (сильно, $\pi^-p\to\Lambda K^0$), а распадаются медленно (слабо).
\end{enumerate}

\paragraph{Смешивание кварков (угол Кабиббо).} $W$ связывает $u$ не с $d$, а с комбинацией
$d'=d\cos\theta_C+s\sin\theta_C$, $\theta_C\approx13^\circ$. Поэтому процессы с изменением странности ($s\to u$) подавлены
множителем $\sin^2\theta_C\approx0{,}05$ по сравнению с $d\to u$. Обобщение на три поколения — матрица
Кабиббо--Кобаяси--Маскавы; её комплексная фаза ответственна за нарушение $CP$ (наблюдается в распадах $K$- и $B$-мезонов).

\begin{figure}[htbp]
\centering
\begin{tikzpicture}[scale=0.85, photon/.style={decorate, decoration={snake, amplitude=1.1mm, segment length=2.6mm}, thick},
  fermion/.style={thick, postaction={decorate}, decoration={markings, mark=at position 0.55 with {\arrow{Stealth[length=2mm]}}}}]
  % распад мюона
  \begin{scope}
   \draw[fermion] (0.5,0) -- (1,1.2); \draw[fermion] (1,1.2) -- (0.5,2.6);
   \draw[photon] (1,1.2) -- (2.3,1.5);
   \draw[fermion] (2.3,1.5) -- (2.6,2.6);
   \draw[fermion] (3.4,2.3) -- (2.3,1.5);
   \node[below, font=\small] at (0.5,0) {$\mu^-$}; \node[above, font=\small] at (0.5,2.6) {$\nu_\mu$};
   \node[above, font=\small] at (2.6,2.6) {$e^-$}; \node[right, font=\small] at (3.4,2.3) {$\bar\nu_e$};
   \node[below, font=\small] at (1.8,1.2) {$W^-$};
   \node[font=\small] at (1.7,-0.7) {(а) $\mu^-\to e^-\bar\nu_e\nu_\mu$};
  \end{scope}
  % распад нейтрона
  \begin{scope}[xshift=5.2cm]
   \draw[fermion] (0.0,0) -- (0.0,2.6); \draw[fermion] (0.3,0) -- (0.3,2.6);
   \draw[fermion] (0.7,0) -- (1.0,1.2); \draw[fermion] (1.0,1.2) -- (0.7,2.6);
   \draw[photon] (1.0,1.2) -- (2.3,1.5);
   \draw[fermion] (2.3,1.5) -- (2.6,2.6);
   \draw[fermion] (3.4,2.3) -- (2.3,1.5);
   \node[below, font=\small] at (0.35,0) {$n\,(u\,d\,d)$}; \node[above, font=\small] at (0.35,2.6) {$p\,(u\,d\,u)$};
   \node[above, font=\small] at (2.6,2.6) {$e^-$}; \node[right, font=\small] at (3.4,2.3) {$\bar\nu_e$};
   \node[below, font=\small] at (1.8,1.2) {$W^-$};
   \node[font=\small] at (1.7,-0.7) {(б) $n\to pe^-\bar\nu_e$};
  \end{scope}
  % распад пиона
  \begin{scope}[xshift=10.4cm]
   \draw[fermion] (0.2,0) -- (0.9,1.2);
   \draw[fermion] (0.9,1.2) -- (1.6,0);
   \draw[photon] (0.9,1.2) -- (1.6,2.1);
   \draw[fermion] (2.6,2.0) -- (1.6,2.1);
   \draw[fermion] (1.6,2.1) -- (1.9,3.0);
   \node[below, font=\small] at (0.2,0) {$u$}; \node[below, font=\small] at (1.6,0) {$\bar d$};
   \node[right, font=\small] at (2.6,2.0) {$\mu^+$}; \node[above, font=\small] at (1.9,3.0) {$\nu_\mu$};
   \node[left, font=\small] at (1.2,1.8) {$W^+$};
   \node[font=\small] at (1.3,-0.7) {(в) $\pi^+\to\mu^+\nu_\mu$};
  \end{scope}
\end{tikzpicture}
\caption{Диаграммы слабых распадов (время — снизу вверх): (а) распад мюона; (б) бета-распад нейтрона на кварковом уровне
($d\to u$, кварки-«наблюдатели» не участвуют); (в) распад $\pi^+$ — аннигиляция $u\bar d$ в виртуальный $W^+$.}
\label{fig:16:decays}
\end{figure}

\subsection{Законы сохранения в слабых процессах}

Сохраняются: энергия, импульс, момент, электрический заряд, барионное число, лептонные числа $L_e$, $L_\mu$, $L_\tau$
(в процессах; нейтринные осцилляции их смешивают). Нарушаются: аромат кварков (странность, очарование и др.),
чётность $P$, зарядовое сопряжение $C$, (слабо) $CP$. Пример запрещённого процесса: $\mu^-\to e^-+\gamma$ (нарушает
$L_e$ и $L_\mu$); экспериментальная доля таких распадов меньше $10^{-12}$.

\section{Математические выводы}

\begin{derivation}[label=der:16:sargent]{закон Сарджента $\Gamma\propto G_F^2E^5$ из размерности}
\emph{Допущения:} энерговыделение $E_0$ (для мюона $E_0\approx m_\mu c^2$) много больше масс продуктов; амплитуда
пропорциональна $G_F$ (низшего порядка), значит ширина $\Gamma\propto G_F^2$.

В единицах $\hbar=c=1$ размерность $[G_F]=\text{энергия}^{-2}$, $[\Gamma]=\text{энергия}$. Единственный оставшийся масштаб — $E_0$.
Следовательно,
\[
  \Gamma=C\,G_F^2E_0^5\qquad(\text{размерность: }E^{-4}\cdot E^5=E).
\]
Безразмерный коэффициент $C$ определяется фазовым объёмом трёх частиц; для распада мюона $C=1/(192\pi^3)$.
Восстанавливая $\hbar$ и $c$: $\Gamma=\dfrac{G_F^2E_0^5}{192\pi^3(\hbar c)^6}$, $\tau=\hbar/\Gamma$.

\emph{Численная проверка для мюона:}
$G_F^2(m_\mu c^2)^5=(1{,}166\cdot10^{-5}\ \text{ГэВ}^{-2})^2\cdot(0{,}10566\ \text{ГэВ})^5=1{,}36\cdot10^{-10}\cdot1{,}316\cdot10^{-5}\ \text{ГэВ}=1{,}79\cdot10^{-15}$~ГэВ;
делим на $192\pi^3=5953$: $\Gamma=3{,}0\cdot10^{-19}$~ГэВ $=3{,}0\cdot10^{-10}$~эВ. Тогда
$\tau=\hbar/\Gamma=6{,}58\cdot10^{-16}\ \text{эВ}\cdot\text{с}/3{,}0\cdot10^{-10}\ \text{эВ}=2{,}2\cdot10^{-6}$~с — совпадает с измеренным временем жизни.
\end{derivation}

\begin{derivation}[label=der:16:tau]{время жизни $\tau$-лептона по универсальности}
Ширина канала $\tau\to e\bar\nu_e\nu_\tau$ вычисляется по той же формуле, что и для мюона, с заменой $m_\mu\to m_\tau$:
$\Gamma(\tau\to e\nu\bar\nu)=\Gamma_\mu(m_\tau/m_\mu)^5$. Если доля этого канала $B_e\approx0{,}178$, то полная ширина
$\Gamma_\tau=\Gamma(\tau\to e\nu\bar\nu)/B_e$, и
\begin{align*}
  \tau_\tau&=\tau_\mu\Bigl(\frac{m_\mu}{m_\tau}\Bigr)^5B_e
  =2{,}197\cdot10^{-6}\ \text{с}\cdot\Bigl(\frac{105{,}66}{1776{,}9}\Bigr)^5\cdot0{,}178\\
  &=2{,}197\cdot10^{-6}\cdot7{,}47\cdot10^{-7}\cdot0{,}178\ \text{с}\approx2{,}9\cdot10^{-13}\ \text{с}.
\end{align*}
Измеренное значение $2{,}90\cdot10^{-13}$~с. Согласие подтверждает лептонную универсальность: $\tau$ и $\mu$ взаимодействуют с
$W$ одинаково. Доля $B_e$ сама объясняется подсчётом каналов: $e\nu\nu$, $\mu\nu\nu$ и $d\bar u$ в трёх цветах — всего
$\approx5$ равноправных каналов, $B_e\approx1/5$.
\end{derivation}

\begin{derivation}[label=der:16:neutron]{оценка времени жизни нейтрона по закону Сарджента}
Сравним с мюоном, считая матричный элемент того же порядка: $\tau_n/\tau_\mu\sim(m_\mu c^2/Q_n)^5$ (с точностью до
множителей фазового объёма, которые для $Q\sim m_ec^2$ отличаются заметно). $(105{,}7/0{,}78)^5\approx(135)^5\approx4{,}5\cdot10^{10}$,
т.\,е. $\tau_n\sim4{,}5\cdot10^{10}\cdot2{,}2\cdot10^{-6}$~с $\sim10^5$~с. Реальное значение 880~с: при $Q\sim m_ec^2$ электрон
нерелятивистский, и закон $Q^5$ лишь грубая оценка; но порядок «очень долго по сравнению с мюоном» объясняется именно малым
энерговыделением. Главный вывод: \emph{время жизни слабых распадов определяется прежде всего энерговыделением}.
\end{derivation}

\begin{derivation}[label=der:16:pion]{кинематика распада $\pi^+\to\mu^+\nu_\mu$}
Пион покоится; продукты разлетаются с равными по модулю импульсами $p$. Сохранение энергии ($m_\nu\approx0$):
\begin{gather*}
  m_\pi c^2=\sqrt{p^2c^2+m_\mu^2c^4}+pc\quad\Longrightarrow\\
  pc=\frac{(m_\pi^2-m_\mu^2)c^2}{2m_\pi}=\frac{139{,}57^2-105{,}66^2}{2\cdot139{,}57}\ \text{МэВ}=29{,}8\ \text{МэВ}.
\end{gather*}
Кинетическая энергия мюона $T_\mu=\sqrt{p^2c^2+m_\mu^2c^4}-m_\mu c^2=4{,}1$~МэВ — моноэнергетический мюон (двухчастичный распад).
Спины: $J_\pi=0$, импульсы противоположны, значит проекции спинов $\mu^+$ и $\nu_\mu$ на свои импульсы одинаковы: левое $\nu_\mu$
требует левого $\mu^+$ — отсюда полная поляризация мюонов от распада пионов.
\end{derivation}

\section{Итоговые формулы}

\keyformula{\frac{G_F}{\sqrt2}=\frac{g^2}{8M_W^2},\qquad \frac{G_F}{(\hbar c)^3}=1{,}166\cdot10^{-5}\ \text{ГэВ}^{-2},\qquad M_Z=\frac{M_W}{\cos\theta_W}}
 {$g$ — константа слабой связи, $\theta_W$ — угол Вайнберга ($\sin^2\theta_W\approx0{,}23$); в первой формуле $\hbar=c=1$.}
 {энергии $\ll M_Wc^2$ (эффективная теория Ферми).}
 {связь константы Ферми с массой переносчика; «слабость» как следствие большой $M_W$.}

\keyformula{\Gamma_\mu=\frac{G_F^2(m_\mu c^2)^5}{192\pi^3(\hbar c)^6},\qquad \Gamma\propto G_F^2E_0^5}
 {$\Gamma=\hbar/\tau$ — ширина, $E_0$ — энерговыделение.}
 {$E_0\gg$ масс продуктов; низший порядок по $G_F$.}
 {время жизни мюона, $\tau$-лептона; оценки времён бета-распадов.}

\keyformula{\mu^-\to e^-\bar\nu_e\nu_\mu,\quad n\to pe^-\bar\nu_e,\quad \pi^+\to\mu^+\nu_\mu,\quad \Lambda\to p\pi^-}
 {$\tau_\mu=2{,}2$~мкс, $\tau_n\approx880$~с, $\tau_{\pi^\pm}=2{,}6\cdot10^{-8}$~с, $\tau_\Lambda=2{,}6\cdot10^{-10}$~с.}
 {сохраняются $Q$, $B$, $L_e$, $L_\mu$; странность меняется на 1 (для $\Lambda$).}
 {типичные слабые распады; проверка законов сохранения.}

\section{Физическая интерпретация}

\textbf{Почему «слабое».} Константа $g$ не мала, но при энергиях $E\ll M_Wc^2$ виртуальный $W$ сильно «вне массовой
поверхности», и амплитуда подавлена множителем $(E/M_Wc^2)^2$. При энергиях LHC слабое взаимодействие по силе
сопоставимо с электромагнитным.

\textbf{Шкала времён жизни.} При одной константе времена слабых распадов различаются на 13 порядков (от
$3\cdot10^{-13}$~с для $\tau$ до 880~с для нейтрона) — из-за степени $E_0^5$ и подавлений (Кабиббо, спиральность).

\textbf{Асимметрия природы.} Нарушение $P$ и $C$ означает, что «зеркальный мир» и «мир из антивещества» по отдельности
физически отличимы от нашего; нарушение $CP$ — одно из условий (Сахаров, 1967) возникновения барионной асимметрии
Вселенной, хотя величины $CP$-нарушения в СМ для её объяснения недостаточно.

\section{Что необходимо знать на экзамене}

\begin{examcheck}
\checkhead{Определения}
\begin{itemize}
 \item слабое взаимодействие: участники, переносчики $W^\pm$, $Z^0$ (массы, спин 1), заряженные и нейтральные токи;
 \item спиральность; нарушение чётности (опыт Ву); угол Кабиббо; лептонные числа.
\end{itemize}
\checkhead{Формулы}
\begin{itemize}
 \item $r_W=\hbar/M_Wc$; $G_F/\sqrt2=g^2/8M_W^2$; $M_Z=M_W/\cos\theta_W$;
 \item $\Gamma\propto G_F^2E_0^5$ (закон Сарджента), $\tau_\tau=\tau_\mu(m_\mu/m_\tau)^5B_e$.
\end{itemize}
\checkhead{Выводы}
\begin{itemize}
 \item закон Сарджента из размерности и численная проверка для мюона;
 \item время жизни $\tau$-лептона по универсальности; кинематика $\pi\to\mu\nu$ и спиральное подавление $\pi\to e\nu$.
\end{itemize}
\checkhead{Интерпретации}
\begin{itemize}
 \item кварковые диаграммы распадов $n$, $\pi^+$, $\Lambda$; почему странные частицы рождаются быстро, а распадаются медленно;
 \item электрослабое объединение и открытие $W$, $Z$, нейтральных токов; три поколения нейтрино по ширине $Z$.
\end{itemize}
\end{examcheck}
